%%%%%%%%%%%%%%%%%%%%%%%%%%%%%%%%%%%%%%%%%%%%%%%%%%%%%%%%%%%%%%%%%%%%%%%%%%%%
% CAPÍTULO 16: BIBLIOGRAFÍA PRELIMINAR
%%%%%%%%%%%%%%%%%%%%%%%%%%%%%%%%%%%%%%%%%%%%%%%%%%%%%%%%%%%%%%%%%%%%%%%%%%%%

\chapter{Bibliografía preliminar}
\label{ch:bibliografia-preliminar}

\section{Criterios de organización}

La bibliografía se organiza en cinco bloques temáticos coherentes con las categorías A--E de la revisión sistemática de la literatura (SLR). Las referencias completas se presentan una sola vez en la bibliografía automática; aquí se identifican las publicaciones clave mediante sus claves y citas (Tabla~\ref{tab:bloques-bibliografia}).

\begin{table}[H]
  \centering
  \begin{tabularx}{\textwidth}{c X l}
    \hline
    \textbf{Bloque} & \textbf{Contenido} & \textbf{Prioridad} \\
    \hline
    A & Fundamentos y arquitectura de Gemelos Digitales & Alta \\
    B & Gemelos Digitales aplicados a movilidad compartida & Muy alta \\
    C & Datos, integración y sincronización operativa & Alta \\
    D & Simulación, comportamiento de usuarios y escenarios & Muy alta \\
    E & \textit{Matching}, rutas, optimización y evaluación & Muy alta \\
    \hline
  \end{tabularx}
  \caption{Criterios de organización de la bibliografía.}
  \label{tab:bloques-bibliografia}
\end{table}

\section{Referencias clave por bloque}

\noindent\textbf{Bloque A: fundamentos y arquitectura de Gemelos Digitales.} \cite{A1,A2,A3,A5,A9}.

\noindent\textbf{Bloque B: Gemelos Digitales y movilidad compartida.} \cite{B1,B4,B5,B6}.

\noindent\textbf{Bloque C: datos, integración y sincronización.} \cite{C1,C3,C8}.

\noindent\textbf{Bloque D: simulación y escenarios.} \cite{D2,D4,D10,D12,D14,D15}.

\noindent\textbf{Bloque E: \textit{matching}, rutas y evaluación.} \cite{E16,E17,E19,E20,E21,E22}.

\section{Metodología de la revisión}

\begin{itemize}
    \item Revisión sistemática de Kitchenham et al. \cite{Kitchenham2009}.
    \item Aplicación de Alfonso et al. \cite{Alfonso2021}.
\end{itemize}

\section{Objetivo bibliográfico}

\begin{itemize}
    \item Base final de trabajo: entre 50 y 77 referencias relevantes del corpus principal, a las que se suman las referencias metodológicas y normativas (Kitchenham, Hevner, Peffers, ISO/IEC 25010, etc.).
    \item Prioridad: calidad y relación directa con el problema, no cantidad.
    \item Criterio: cada referencia debe poder conectarse con al menos una decisión de diseño, un requisito o una métrica del TFM.
\end{itemize}
